\documentclass[a4paper]{article}
\usepackage{amsfonts}
\usepackage{amsmath}
\usepackage{amssymb}
\usepackage{graphicx}

\begin{document}
  
\noindent \large Auteur: Michael Zielinski \\
\noindent \large Studierichting: Informatica\\
\noindent \large Oefeningenreeks 7A nr. 2.1\\

\medskip

\normalsize

Bereken: \\

$ \lim\limits_{(x,y,z)\rightarrow(1,1,1)}
\dfrac{x^2+y^2+z^2}{1-x-y-z} $ \\

Omdat de noemer in $(1,1,1)$ niet nul wordt, mogen we rechtstreeks invullen. \\

$ 1-x-y-z = 1-1-1-1 $ \\

$ = -2 $ \\

$ x^2+y^2+z^2 = 1^2+1^2+1^2 $ \\

$ = 3 $ \\

Dus: \\

$ \lim\limits_{(x,y,z)\rightarrow(1,1,1)}
\dfrac{x^2+y^2+z^2}{1-x-y-z}
= \dfrac{3}{-2} $ \\

$ = -\dfrac{3}{2} $ \\

\begin{thebibliography}{999}
%\bibitem{a} Referentie
\end{thebibliography}

\end{document}